
\documentclass[conference]{IEEEtran}

\usepackage{cite}
\usepackage{amsmath,amssymb,amsfonts}
\usepackage{graphicx}
\usepackage{booktabs}
\usepackage{array}
\usepackage{multirow}
\usepackage{url}
\usepackage{hyperref}

\begin{document}

\title{Keystroke Dynamics for Emotion Recognition: A Survey of Methods, Applications, and Future Directions}

\author{
    \IEEEauthorblockN{Author Name}
    \IEEEauthorblockA{
        Department Name, Institution Name\\
        email@example.com
    }
}

\maketitle

\begin{abstract}
This paper provides a comprehensive review of keystroke dynamics as a behavioral biometric for emotion recognition. It synthesizes recent advances in datasets, feature extraction, machine learning models, and application domains such as cybersecurity, human-computer interaction, and mental health monitoring. Key challenges include individual variability, contextual influences, and the need for robust, non-intrusive methods. The survey highlights multimodal approaches integrating keystroke patterns with physiological signals and text semantics, while addressing limitations in existing studies. Emphasis is placed on scalable solutions for real-world deployment, including adaptive systems that respond to emotional states. Future directions include improving personalization, enhancing interpretability, and expanding datasets to diverse populations. This work serves as a roadmap for researchers aiming to advance emotion-aware computing through keystroke-based analytics.
\end{abstract}

\begin{IEEEkeywords}
['Keystroke Dynamics', 'Emotion Recognition', 'Behavioral Biometrics', 'Machine Learning', 'Human-Computer Interaction', 'Multimodal Systems']
\end{IEEEkeywords}


\section{Introduction}

The integration of emotion recognition into human-computer interaction (HCI) has gained significant traction as adaptive systems become increasingly prevalent in modern computing environments. Traditional approaches to emotion detection have relied heavily on physiological sensors, facial expressions, speech analysis, and input device data, often requiring intrusive hardware or specialized equipment. These methods, while effective in controlled settings, face challenges related to cost, user privacy, and real-time applicability. Keystroke dynamics, a behavioral biometric derived from typing patterns, has emerged as a promising non-intrusive alternative. By analyzing features such as dwell time, flight time, latency, and session duration, keystroke dynamics offers a cost-effective means of capturing emotional states without requiring physical sensors. This review synthesizes existing research on keystroke dynamics for emotion recognition, highlighting its potential, limitations, and future directions.  

The role of keystroke dynamics in emotion recognition is rooted in its ability to capture subtle variations in human behavior influenced by emotional states. Studies have demonstrated that emotional arousal and valence significantly impact typing patterns, with negative emotions often associated with reduced typing speed and increased latency, while positive emotions may correlate with faster keystrokes and shorter dwell times. These patterns, when combined with machine learning techniques such as support vector machines (SVM), fuzzy logic, and convolutional neural networks (CNN), enable the classification of emotional states with varying degrees of accuracy. For instance, research has shown that short typing windows (30 seconds) can achieve performance comparable to longer sessions, suggesting the feasibility of real-time emotion detection. However, the reliability of these models depends on the quality and diversity of datasets, as well as the robustness of feature extraction methods.  

Existing literature on keystroke dynamics for emotion recognition reveals both progress and gaps. Early studies focused on static datasets and limited emotional categories, often relying on self-reported labels that introduce potential bias. More recent work has addressed these limitations by incorporating multimodal approaches, integrating keystroke data with physiological signals like heart rate variability (HRV) and text analysis using natural language processing (NLP) tools such as NLTK and JAWS. These hybrid models have demonstrated improved accuracy, particularly in distinguishing between positive and negative emotional states. For example, one study found that NLTK outperformed JAWS in text-based emotion detection, while HRV correlated strongly with emotional states across multiple datasets. However, the integration of physiological signals introduces challenges related to noise, signal variability, and the need for calibrated hardware, which may limit scalability in real-world applications.  

The methodology of keystroke dynamics research has evolved from isolated feature extraction to systematic analysis of emotional variance. Systematic reviews of articles from 2013 to the present have identified trends in dataset development, algorithmic advancements, and application domains. For instance, benchmark datasets such as the IAPS (International Affective Picture System) and SAM (Self-Assessment Manikin) have been widely used to evaluate emotional states, while newer studies emphasize the need for culturally diverse and context-aware datasets. The use of machine learning frameworks, including long short-term memory (LSTM) networks and conditional diffusion models, has enabled more nuanced modeling of emotional dynamics, particularly in scenarios requiring continuous monitoring. However, the scalability of these models remains a challenge, as demonstrated by the need for large-scale datasets (e.g., 100,000+ users) to achieve reliable performance metrics.  

Despite its potential, keystroke dynamics for emotion recognition faces several limitations. One major challenge is the variability of individual typing patterns, which can obscure emotional signals and reduce model generalizability. Additionally, external factors such as device type, keyboard layout, and user familiarity with the interface can introduce confounding variables. For example, studies using smartphone typing data have shown that touch-based interactions differ significantly from desktop keyboards, necessitating tailored feature extraction methods. Furthermore, the reliance on self-reported labels for emotional states raises concerns about data validity, as subjective interpretations may not align with objective behavioral metrics. These limitations underscore the need for more robust validation frameworks and the integration of multimodal data to mitigate biases and improve accuracy.  

The implications of keystroke dynamics extend beyond emotion recognition to broader applications in cybersecurity, mental health monitoring, and adaptive computing systems. In cybersecurity, keystroke dynamics has been leveraged for non-intrusive user authentication, with recent advancements such as TypeNet demonstrating scalability for large-scale internet authentication systems. In mental health, the ability to detect mood disturbances through typing metadata has opened avenues for early intervention and personalized care. For instance, the DeepMood framework achieved 90.31\% accuracy in depression prediction using session-level typing data, highlighting the potential of behavioral biometrics in healthcare. However, these applications require careful consideration of ethical and privacy concerns, particularly in contexts involving sensitive user data.  

Future research directions in keystroke dynamics for emotion recognition include addressing the limitations of unimodal approaches by integrating multimodal data, improving the generalizability of models through diverse datasets, and enhancing real-time processing capabilities. The development of frameworks that account for cultural and contextual variability, such as the proposed human micro-expression-based social behavioral biometrics (SBB), could further expand the applicability of keystroke dynamics. Additionally, the integration of affective computing with ambient assisted living systems offers opportunities for continuous, non-intrusive monitoring of emotional states in aging populations and individuals with neurodegenerative diseases.  

In conclusion, keystroke dynamics represents a transformative approach to emotion recognition, offering a non-intrusive, cost-effective alternative to traditional methods. While challenges such as individual variability, data volume, and external interference persist, the integration of advanced machine learning techniques and multimodal data has significantly enhanced its feasibility. As research continues to address these limitations, the potential of keystroke dynamics to revolutionize HCI, cybersecurity, and mental health monitoring becomes increasingly evident. Future work should prioritize the development of scalable, context-aware systems that balance accuracy with user privacy, ensuring the broad adoption of this promising technology.

\section{Literature Review}

The field of keystroke dynamics has evolved significantly as a behavioral biometric, offering non-intrusive methods for user authentication, emotion recognition, and mental state monitoring. Early research focused on its role in cybersecurity, where keystroke features such as dwell time and flight time were leveraged to distinguish users and detect anomalies. However, recent advancements have expanded its applications to affective computing, where keystroke dynamics are increasingly viewed as a proxy for emotional and cognitive states. This synthesis highlights the trajectory of keystroke dynamics research, its integration with multimodal approaches, and the challenges that remain in achieving robust, scalable systems for emotion recognition and user authentication.  

The foundational work on keystroke dynamics in cybersecurity established its potential as a cost-effective, non-intrusive alternative to traditional biometric methods. Studies from 2013 onward emphasized the importance of benchmark datasets and algorithmic advancements, such as the use of deep learning models like convolutional neural networks (CNNs) and support vector machines (SVMs) to improve classification accuracy. These efforts demonstrated that keystroke dynamics could achieve performance comparable to other biometric modalities, particularly in scenarios where physical sensors are impractical. However, the transition from authentication to emotion recognition introduced new complexities, as the relationship between typing patterns and emotional states is inherently more nuanced. Early explorations in this domain highlighted the need for controlled experiments and the integration of physiological signals, such as heart rate variability, to complement keystroke data.  

The shift toward emotion recognition in human-computer interaction (HCI) has been driven by the desire to create adaptive systems that respond to users' emotional states. Research in this area has revealed that keystroke dynamics can capture subtle variations in typing behavior influenced by emotions. For instance, studies have shown that emotional states such as anger or sadness can alter typing speed, latency, and accuracy, with negative emotions often associated with reduced performance. These findings have been corroborated through controlled experiments using standardized stimuli like the International Affective Picture System (IAPS) and self-assessment manikin (SAM) ratings. However, the reliability of keystroke-based emotion recognition remains contested, as individual variability and contextual factors--such as cultural differences or task complexity--introduce significant noise into the data. This has led to calls for multimodal approaches that combine keystroke dynamics with other modalities, such as text analysis and physiological signals, to enhance accuracy and robustness.  

A key challenge in keystroke dynamics research is the underexplored potential of real-time emotion detection. While existing studies have demonstrated the feasibility of classifying emotional states using features like session time and latency, the practical application of these methods in dynamic environments remains limited. For example, short writing windows (30 seconds) have been shown to achieve similar accuracy to longer sessions, suggesting that real-time analysis is possible. However, the integration of keystroke data into continuous emotion recognition systems requires addressing issues such as data volume, external interference, and the need for large, diverse datasets. Recent work has attempted to mitigate these challenges by constructing specialized datasets and employing advanced machine learning techniques, including fuzzy logic and ensemble methods, to model complex emotional patterns.  

The intersection of keystroke dynamics with affective computing has also sparked interest in applications beyond traditional authentication. For instance, the ability to detect stress or mood disturbances through typing patterns has been explored in the context of mental health monitoring and cyberbullying detection. Studies have demonstrated that metrics such as typing pressure and mouse contact differences can serve as indicators of stress levels, with implications for developing non-intrusive health monitoring systems. Similarly, the use of keystroke dynamics in social media contexts has been proposed as a means to identify emotional states in user-generated content, offering potential applications in personalized computing and mental health interventions. However, these applications are still in their nascent stages, constrained by the need for more granular data and the limitations of unimodal approaches.  

Despite these advancements, several gaps persist in the literature. The majority of studies focus on predefined emotional states, often relying on self-reported labels that introduce bias. Additionally, the generalizability of keystroke-based emotion recognition across diverse populations remains unclear, as factors such as age, gender, and cultural background are known to influence typing behavior. Furthermore, the integration of keystroke dynamics with other biometric modalities--such as facial expressions or voice analysis--has been underexplored, despite the potential for synergistic effects. The lack of standardized evaluation metrics and benchmark datasets also hinders progress, as researchers often use disparate methodologies and datasets, making it difficult to compare results across studies.  

Looking ahead, the future of keystroke dynamics research lies in addressing these limitations through interdisciplinary collaboration and methodological innovation. The development of larger, more representative datasets, coupled with advanced machine learning frameworks, could enhance the accuracy and scalability of emotion recognition systems. Additionally, the exploration of contextual factors--such as task relevance and environmental conditions--may provide deeper insights into how emotional states influence typing patterns. The integration of keystroke dynamics with emerging technologies, such as wearable sensors and ambient intelligence systems, could further expand its applications in healthcare, cybersecurity, and human-computer interaction. As the field continues to evolve, the emphasis on non-intrusive, real-time emotion sensing will likely drive new paradigms in affective computing, positioning keystroke dynamics as a cornerstone of next-generation intelligent systems.

\section{Methodology Analysis}

The methodology analysis of keystroke dynamics in emotion recognition, user authentication, and affective computing reveals a diverse array of approaches, datasets, and algorithmic strategies aimed at leveraging typing patterns for non-intrusive behavioral biometrics. Central to these studies is the extraction of temporal and spatial features from keystroke data, such as dwell time (the duration a key is pressed), flight time (the interval between key presses), and latency (the time between consecutive keystrokes). These features are often combined with statistical models, machine learning algorithms, and multimodal integration to enhance the accuracy and robustness of emotion detection, authentication, and user behavior analysis. The systematic review of methodologies across multiple studies highlights both the advancements in computational techniques and the persistent challenges in achieving reliable, scalable, and culturally adaptive systems.  

A recurring theme in the methodology is the use of supervised machine learning frameworks to classify emotional states or authenticate users based on keystroke dynamics. Algorithms such as support vector machines (SVM), convolutional neural networks (CNN), and long short-term memory (LSTM) networks have been employed to model the relationship between typing patterns and emotional or behavioral states. For instance, SVM has been used to distinguish between emotional states by analyzing features like keystroke duration and latency, while CNNs have been applied to time-windowed keystroke data to capture spatial-temporal correlations. LSTM architectures, with their ability to process sequential data, have been particularly effective in modeling the dynamic nature of typing patterns, enabling the detection of subtle variations that correlate with emotional or cognitive states. These models are often trained on benchmark datasets, such as those derived from standardized emotion elicitation protocols (e.g., IAPS stimuli) or real-world user interactions, to ensure generalizability.  

The integration of keystroke dynamics with other modalities, such as physiological signals (e.g., heart rate variability via PPG) and text analysis (e.g., sentiment classification using NLTK or JAWS), has emerged as a critical strategy to overcome the limitations of unimodal approaches. Multimodal systems combine features from multiple sources to improve detection accuracy, as demonstrated in studies where heart rate variation correlated with emotional states across datasets. Similarly, text analysis tools have been used to complement keystroke features, with natural language processing (NLP) techniques like sentiment lexicons and topic modeling providing additional context for emotion recognition. However, the reliance on self-reported labels or subjective annotations introduces potential biases, necessitating the development of more objective evaluation metrics and larger, diverse datasets to mitigate cultural and individual variability.  

The methodology also emphasizes the importance of experimental design in capturing the nuances of keystroke dynamics. Controlled experiments, such as those using standardized emotion induction protocols (e.g., SAM ratings with IADS-2 stimuli), have been employed to isolate the effects of emotional states on typing patterns. These studies often involve statistical analyses, including two-way ANOVA and cross-validation, to assess the significance of features like duration, latency, and accuracy in different emotional contexts. For example, significant arousal effects on keystroke duration and latency have been observed, with minimal impact on accuracy, suggesting that emotional states primarily influence the speed and timing of typing rather than the correctness of input. Such findings underscore the need for adaptive models that account for individual differences and contextual factors, such as stress or cognitive load, which can modulate keystroke patterns.  

Despite these advancements, several methodological challenges persist. The scalability of keystroke-based systems remains a critical issue, particularly in large-scale applications like internet authentication or mental health monitoring. Studies have highlighted the trade-offs between different loss functions and dataset configurations, with some approaches prioritizing accuracy over computational efficiency. Additionally, the volume and complexity of keystroke data pose challenges in real-time processing, requiring optimized feature extraction and model architectures. The integration of keystroke dynamics with other behavioral biometrics, such as mouse movements or touchscreen interactions, has also been explored to enhance the robustness of emotion recognition systems, though this introduces additional layers of complexity in data fusion and interpretation.  

The review further underscores the role of systematic literature reviews in identifying gaps and trends in keystroke dynamics research. These reviews have analyzed the strengths and weaknesses of existing methodologies, such as the underexplored potential of keystroke dynamics in real-time emotion detection compared to physiological or visual methods. They have also emphasized the need for standardized evaluation frameworks to compare the performance of different algorithms and datasets. For instance, the construction of new benchmark datasets, such as those focused on depression prediction or stress detection, has been proposed to address the limitations of existing data sources and improve the generalizability of findings.  

In conclusion, the methodology analysis of keystroke dynamics reveals a convergence of computational techniques, experimental rigor, and interdisciplinary approaches to advance non-intrusive emotion recognition and authentication. While the integration of multimodal data and the development of scalable models have demonstrated the feasibility of these systems, challenges such as data variability, cultural adaptation, and real-time processing requirements necessitate further innovation. Future research should prioritize the creation of larger, diverse datasets, the refinement of hybrid models that combine keystroke dynamics with other behavioral cues, and the exploration of novel algorithms to enhance the accuracy and applicability of these technologies across domains like healthcare, cybersecurity, and human-computer interaction.

\section{Challenges and Future Directions}

The integration of keystroke dynamics into emotion recognition, user authentication, and affective computing represents a significant advancement in non-intrusive behavioral biometrics. However, the field remains at a critical juncture, where promising developments coexist with persistent challenges that necessitate further exploration. Central to these challenges is the inherent variability in human behavior, which complicates the extraction of consistent and reliable features from keystroke data. While existing methodologies, such as the use of dwell time, flight time, and session-based analysis, have demonstrated potential for detecting emotional states or authenticating users, their effectiveness is often constrained by factors such as individual differences, contextual influences, and the limitations of unimodal approaches. For instance, studies have shown that emotional states can significantly alter typing patterns, with negative emotions often associated with reduced speed and increased latency, while positive states may enhance typing efficiency. However, these patterns are not universally consistent, as cultural, demographic, and situational variables introduce noise that undermines the reliability of keystroke-based models. This variability underscores the need for more robust frameworks that account for such factors, particularly in real-world applications where user behavior is influenced by a multitude of external and internal stimuli.

A key limitation in current research is the reliance on self-reported labels for emotional states, which introduces potential biases and inaccuracies. While self-reported data provides a direct measure of subjective experience, it may not align with objective behavioral indicators, leading to discrepancies in model training and evaluation. Additionally, the scarcity of large, diverse, and annotated datasets for keystroke dynamics research hampers the development of generalized models. Most studies have focused on controlled environments, often using limited datasets that may not capture the complexity of real-world interactions. For example, while some experiments have demonstrated high accuracy in classifying emotional states using features like keystroke duration and latency, these results are frequently validated in isolated settings, raising questions about their scalability and applicability to broader populations. The challenge of data volume and quality is further compounded by the need for continuous, real-time analysis, which requires efficient algorithms capable of processing dynamic and heterogeneous input streams without compromising performance.

Another critical issue is the integration of keystroke dynamics with other modalities to enhance accuracy and robustness. While unimodal approaches have shown promise, their limitations in capturing the full spectrum of human behavior highlight the necessity of multimodal integration. For instance, combining keystroke data with physiological signals, such as heart rate variability or electrodermal activity, or with text-based features, could provide a more comprehensive understanding of emotional states. However, such integration is not without its own challenges. The fusion of heterogeneous data sources requires sophisticated techniques to align and synchronize features across different modalities, while also addressing issues like signal noise, data ambiguity, and computational complexity. Moreover, the ethical and privacy implications of collecting and analyzing multimodal data, particularly in sensitive contexts like mental health monitoring or cybersecurity, demand careful consideration to ensure user trust and compliance with regulatory standards.

Looking ahead, the future of keystroke dynamics research lies in addressing these challenges through innovative methodologies and interdisciplinary collaboration. One promising direction is the development of adaptive models that can dynamically adjust to individual user behavior and contextual factors. For example, machine learning frameworks such as deep neural networks and fuzzy logic systems could be refined to better capture the nuances of keystroke patterns, while also incorporating contextual metadata to improve prediction accuracy. Additionally, the creation of standardized, publicly available datasets that encompass diverse populations and real-world scenarios would be instrumental in advancing the field. Such datasets could include annotated keystroke data from varied emotional states, demographic groups, and environmental conditions, enabling researchers to train and validate models under more representative conditions. Furthermore, the exploration of edge computing and lightweight algorithms could facilitate real-time processing of keystroke data, making non-intrusive emotion recognition and authentication more feasible for applications such as mobile devices, wearable technology, and remote healthcare monitoring.

The potential applications of keystroke dynamics extend beyond traditional cybersecurity and affective computing domains. In mental health, for instance, the ability to detect subtle changes in typing patterns could enable early identification of mood disturbances or stress-related conditions, offering a non-invasive tool for continuous monitoring. Similarly, in educational and workplace settings, keystroke-based systems could provide insights into user engagement, cognitive load, or emotional well-being, supporting adaptive learning environments and productivity optimization. However, the successful translation of these applications into practical solutions requires overcoming technical, ethical, and usability barriers. For example, ensuring the privacy and security of keystroke data, particularly in scenarios involving sensitive information, is paramount. Additionally, user acceptance and adoption of such systems depend on their perceived utility and ease of integration into existing workflows, which necessitates user-centered design approaches and rigorous validation studies.

In conclusion, while keystroke dynamics has demonstrated significant potential as a non-intrusive method for emotion recognition and authentication, its full realization hinges on addressing existing challenges through interdisciplinary research and innovation. The development of robust, scalable, and ethically sound frameworks will be critical in unlocking the broader applications of this technology. By prioritizing the integration of multimodal data, improving data quality and diversity, and refining algorithmic approaches, the field can move closer to achieving reliable, real-world solutions that enhance human-computer interaction and contribute to advancements in cybersecurity, mental health, and affective computing. The journey toward this goal will require sustained collaboration across disciplines, as well as a commitment to addressing the complex interplay between human behavior, technological capabilities, and societal needs.

\section{Conclusion}

The integration of keystroke dynamics into affective computing, cybersecurity, and human-computer interaction represents a significant advancement in non-intrusive behavioral biometrics. This review synthesizes findings from diverse studies that explore the potential of keystroke dynamics as a tool for emotion recognition, user authentication, and mental state monitoring. While the methodologies and applications vary across these domains, a common thread emerges: the ability of keystroke patterns to capture subtle physiological and psychological cues, offering a cost-effective and scalable alternative to traditional intrusive methods. However, the research also underscores critical challenges that must be addressed to fully realize the potential of keystroke dynamics in real-world scenarios.  

Keystroke dynamics, which encompasses metrics such as dwell time, flight time, latency, and typing speed, has been increasingly leveraged to infer emotional states, cognitive load, and behavioral intent. Studies have demonstrated that emotional fluctuations--whether driven by stress, anxiety, or positive engagement--manifest in distinct typing patterns. For instance, negative emotional states often correlate with reduced typing speed and increased latency, while positive states may exhibit faster, more fluid keystrokes. These observations align with broader theories in affective computing, which posit that human behavior is deeply intertwined with emotional and cognitive processes. The ability to decode such patterns through keystroke analysis opens new avenues for applications in mental health monitoring, adaptive user interfaces, and cybersecurity. However, the reliability of these inferences depends on the granularity of the features extracted and the robustness of the classification models employed.  

A key strength of keystroke dynamics lies in its non-intrusive nature, which makes it particularly appealing for scenarios where physical sensors or user cooperation are impractical. Unlike physiological methods that require wearable devices or invasive data collection, keystroke analysis relies on existing input devices, such as keyboards and touchscreens, making it accessible in a wide range of contexts. This advantage is further amplified by the scalability of the approach, as demonstrated by studies that have successfully applied machine learning algorithms--including support vector machines (SVM), convolutional neural networks (CNN), and long short-term memory (LSTM) networks--to classify emotional states or authenticate users. For example, one study reported a 95\% accuracy rate in distinguishing liars from truth-tellers using keystroke features combined with unexpected questions, while another achieved 90.31\% depression prediction accuracy through session-level typing metadata. These results highlight the potential of keystroke dynamics to serve as a covert yet effective tool for deception detection and mental health assessment.  

Despite these promising outcomes, the research also identifies several limitations that hinder the widespread adoption of keystroke dynamics. One major challenge is the variability in user behavior, which can obscure the relationship between keystroke patterns and emotional states. Individual differences in typing habits, cultural factors, and contextual influences--such as the urgency of a task or the presence of distractions--can introduce noise into the data, reducing the accuracy of classifications. Additionally, the reliance on self-reported labels for emotional states raises concerns about bias, as subjective interpretations may not align with objective behavioral metrics. For instance, studies have noted that while keystroke features can capture arousal levels, they often fail to account for the nuanced interplay between valence (positive/negative emotions) and other psychological dimensions. These limitations underscore the need for multimodal approaches that integrate keystroke data with other signals, such as physiological measurements or text analysis, to enhance the robustness of emotion recognition systems.  

The application of keystroke dynamics in cybersecurity further illustrates its versatility and potential. Traditional authentication methods, such as passwords and biometric scans, are vulnerable to theft or spoofing, whereas keystroke dynamics offers a dynamic, behavior-based alternative. By analyzing the unique typing patterns of users, systems can continuously verify identity without requiring explicit input, making it particularly useful for securing sensitive applications. However, the scalability of these systems remains a critical consideration. While some studies have demonstrated the feasibility of deploying keystroke-based authentication on large datasets, challenges such as data volume, computational overhead, and cross-device compatibility must be addressed to ensure practical implementation. For example, the TypeNet framework, which employs LSTM architectures and cross-device dataset analysis, has shown promising scalability results, but further refinement is needed to optimize performance across diverse user populations and device ecosystems.  

The broader implications of keystroke dynamics extend beyond technical applications, influencing fields such as human-computer interaction (HCI) and ambient assisted living. In HCI, the integration of emotional state detection through keystroke analysis could enable adaptive interfaces that respond to user needs in real time. For instance, systems could adjust task difficulty or provide feedback based on inferred stress levels, enhancing user experience and productivity. In healthcare, keystroke dynamics has been proposed as a tool for monitoring mental health conditions, such as depression or anxiety, by analyzing changes in typing patterns over time. However, the ethical and practical considerations of such applications--particularly in terms of data privacy and the potential for misinterpretation--must be carefully navigated.  

Looking ahead, the research highlights several directions for future work. First, the development of larger, more diverse datasets is essential to improve the generalizability of keystroke-based models. Current studies often rely on limited sample sizes or specific demographic groups, which may not capture the full spectrum of human behavior. Second, the integration of keystroke dynamics with other modalities--such as facial expressions, speech, or physiological signals--could enhance the accuracy and reliability of emotion recognition systems. Third, addressing the cultural and contextual variability in keystroke patterns is crucial for ensuring equitable performance across different populations. Finally, the exploration of real-time applications, such as dynamic authentication or adaptive interfaces, requires further investigation into the computational efficiency and latency of keystroke-based systems.  

In conclusion, the synthesis of findings from this review underscores the transformative potential of keystroke dynamics in affective computing, cybersecurity, and human-computer interaction. While the field has made significant strides in demonstrating the feasibility of non-intrusive emotion recognition and authentication, ongoing research is needed to overcome existing challenges and expand the scope of applications. By addressing limitations in data variability, model robustness, and real-world adaptability, keystroke dynamics can continue to evolve as a cornerstone of behavioral biometrics, offering innovative solutions to complex problems in technology and society. The interdisciplinary nature of this research--spanning psychology, computer science, and engineering--positions keystroke dynamics as a promising avenue for advancing our understanding of human behavior and its integration with intelligent systems.

\begin{thebibliography}{99}

\bibitem{ref1}
Rashik Shadman, Ahmed Anu Wahab, Michael Manno, Matthew Lukaszewski, Daqing Hou, Faraz Hussain, "Keystroke Dynamics: Concepts, Techniques, and Applications," ACM Comput. Surv., 2025. doi: 10.1145/3733103.

\bibitem{ref2}
Stefano Marrone, Carlo Sansone, "Identifying Users' Emotional States through Keystroke Dynamics," 3rd International Conference on Deep Learning Theory and Applications (DeLTA 2022), 2022. doi: 10.5220/0011367300003277.

\bibitem{ref3}
Rinky Solanki, Pragya Shukla, "Estimation of the User's Emotional State by Keystroke Dynamics," International Journal of Computer Applications, 2014.

\bibitem{ref4}
ENYINNAH, CHIDIEBERE, "AN IMPROVED MODEL FOR EMOTION-AWARE COMPUTING DEVICES USING KEYSTROKE DYNAMICS AND SUPPORT VECTOR MACHINE," 2nd International Conference on Education and Development, 2019.

\bibitem{ref5}
Po-Ming Lee, Wei-Hsuan Tsui, Tzu-Chien Hsiao, "The influence of emotion on keyboard typing: an experimental study using visual stimuli," BioMedical Engineering OnLine, 2014. doi: 10.1186/2194-0401-13-81.

\bibitem{ref6}
Anil Kumar K M, Kiran B R, Shreyas B R, Sylvester J Victor, "A Multimodal Approach To Detect User's Emotion," Procedia Computer Science, 2015. doi: 10.1016/j.procs.2015.10.096.

\bibitem{ref7}
Preeti Khanna, M.Sasikumar, "Recognising Emotions from Keyboard Stroke Pattern," International Journal of Computer Applications, 2010.

\bibitem{ref8}
Aicha Maalej, Ilhem Kallel, "Does Keystroke Dynamics tell us about Emotions? A Systematic Literature Review and Dataset Construction," IEEE, 2020.

\bibitem{ref9}
Surjya Ghosh, Niloy Ganguly, Bivas Mitra, Pradipta De, "Evaluating Effectiveness of Smartphone Typing as an Indicator of User Emotion," 7th International Conference on Affective Computing and Intelligent Interaction (ACII), 2017.

\bibitem{ref10}
Rayhan Shikder, Sydur Rahaman, Farzia Afroze, A. B. M. Alim Al Islam, "Keystroke/Mouse Usage Based Emotion Detection and User Identification," NSysS 2017, 2017.

\bibitem{ref11}
Jinfeng Gao, Gangqiang Li, Ruxian Yao, Qiang Liu, Junming Zhang, "Progressive Conditional Diffusion Model for Multistage Spectral Restoration of Remote Sensing Image," IEEE JOURNAL OF SELECTED TOPICS IN APPLIED EARTH OBSERVATIONS AND REMOTE SENSING, 2025. doi: 10.1109/JSTARS.2025.3586437.

\bibitem{ref12}
Po-Ming Lee, Wei-Hsuan Tsui, Tzu-Chien Hsiao, "The Influence of Emotion on Keyboard Typing: An Experimental Study Using Auditory Stimuli," PLOS ONE, 2015. doi: 10.1371/journal.pone.0129056.

\bibitem{ref13}
Nico H. Frijda, "The Laws of Emotion," American Psychologist, 1988.

\bibitem{ref14}
Chieh-Yang Huang, Tristan Labetoulle, Ting-Hao (Kenneth) Huang, Yi-Pei Chen, Hung-Chen Chen, Vallari Srivastava, Lun-Wei Ku, "MoodSwipe: A Soft Keyboard that Suggests Messages Based on User-Specified Emotions," 2017. doi: arXiv:1707.07191v1.

\bibitem{ref15}
Lisa Feldman Barrett, Batja Mesquita, Kevin N. Ochsner, James J. Gross, "The Experience of Emotion," Annu Rev Psychol., 2007.

\bibitem{ref16}
Alejandro Acien, Aythami Morales, John V. Monaco, Ruben Vera-Rodriguez, Julian Fierrez, "TypeNet: Deep Learning Keystroke Biometrics," Journal of LaTeX Class Files, 2021.

\bibitem{ref17}
Merylin Monaro, Chiara Galante, Riccardo Spolaor, Qian Qian Li, Luciano Gamberini, Mauro Conti, Giuseppe Sartori, "Covert lie detection using keyboard dynamics," Scientific Reports, 2018. doi: 10.1038/s41598-018-20462-6.

\bibitem{ref18}
Madiha Tahir, Zahid Halim, Atta Ur Rahman, Muhammad Waqas, Shanshan Tu, Sheng Chen, Zhu Han, "Non-Acted Text and Keystrokes Database and Learning Methods to Recognize Emotions," ACM Trans. Multimedia Comput. Commun. Appl., 2022. doi: 10.1145/3480968.

\bibitem{ref19}
Po-Ming Lee, Wei-Hsuan Tsui, Tzu-Chien Hsiao, "The Influence of Emotion on Keyboard Typing: An Experimental Study Using Auditory Stimuli," PLOS ONE, 2015. doi: 10.1371/journal.pone.0129056.

\bibitem{ref20}
Bokai Cao, Lei Zheng, Chenwei Zhang, Philip S. Yu, Andrea Piscitello, John Zulueta, Olu Ajilore, Kelly Ryan, Alex D. Leow, "DeepMood: Modeling Mobile Phone Typing Dynamics for Mood Detection," KDD'17, 2017. doi: 10.1145/3097983.3098086.

\bibitem{ref21}
Clayton Epp, Michael Lippold, Regan L. Mandryk, "Identifying Emotional States using Keystroke Dynamics," CHI 2011, 2011.

\bibitem{ref22}
Stefanos Xefteris, Nikos Doulamis, Vassiliki Andronikou, Theodora Varvarigou, George Cambourakis, "Behavioral Biometrics in Assisted Living: A Methodology for Emotion Recognition," Engineering, Technology \& Applied Science Research, 2016.

\bibitem{ref23}
Matthias Trojahn, Florian Arndt, Markus Weinmann, Frank Ortmeier, "Emotion Recognition through Keystroke Dynamics on Touchscreen Keyboards," ICEIS-2013, 2013. doi: 10.5220/0004415500310037.

\bibitem{ref24}
A. Ko!akowska, "A Review of Emotion Recognition Methods Based on Keystroke Dynamics and Mouse Movements," IEEE, 2013. doi: 978-1-4673-5637-4.

\bibitem{ref25}
ZAMAN WAHID, ASM HOSSAIN BARI, FAHIM ANZUM, MARINA L. GAVRILOVA, "Human Micro-Expression: A Novel Social Behavioral Biometric for Person Identification," IEEE Access, 2023. doi: 10.1109/ACCESS.2023.3283932.

\bibitem{ref26}
Agata Koakowska, Agnieszka Landowska, "Keystroke Dynamics Patterns While Writing Positive and Negative Opinions," Sensors, 2021. doi: 10.3390/s21175963.

\bibitem{ref27}
Javier Hernandez, Pablo Paredes, Asta Roseway, Mary Czerwinski, "Under Pressure: Sensing Stress of Computer Users," CHI 2014, 2014. doi: 10.1145/2556288.2557165.

\bibitem{ref28}
Wei-Hsuan Tsui, Poming Lee, Tzu-Chien Hsiao, "The Effect of Emotion on Keystroke: An Experimental Study Using Facial Feedback Hypothesis," 35th Annual International Conference of the IEEE EMBS, 2013.

\bibitem{ref29}
LIYING YANG, SHENG-FENG QIN, "A Review of Emotion Recognition Methods From Keystroke, Mouse, and Touchscreen Dynamics," IEEE Access, 2021. doi: 10.1109/ACCESS.2021.3132233.

\bibitem{ref30}
Liying Yang, Shengfeng Qin, "Identify user features impacting keystroke, mouse and touchscreen dynamics," Multimedia Tools and Applications, 2025. doi: 10.1007/s11042-025-21043-2.

\end{thebibliography}

\end{document}
